
\chapter{Methodology}
\label{ch:methodology}

\section{Design principle}

The study is organised so that a negative result is as reportable as a positive one.
Four hypotheses were written down before the experiments they describe, each with a
direction of effect and a decision rule. Where the data contradicted our expectation,
the hypothesis was recorded as not supported rather than reframed after the fact.

Two rules govern the numbers in this report. Every reported value is produced by a
script from a stored artifact, and `scripts/build\_tables.py --verify` fails if a
document quotes a figure the registry does not contain. Test data is touched only
where stated, and the number of times it has been scored is disclosed.

\section{Architectures}

Four models spanning three orders of magnitude in capacity are compared. All are
ImageNet-pretrained except the first, which exists as a deliberate low-capacity
control.

\begin{table}[H]
\centering
\caption{Compared architectures}
\label{tab:arch}
\tabfont
\begin{tabularx}{\textwidth}{@{}lrrY@{}}
\toprule
\textbf{Model} & \textbf{Parameters} & \textbf{Input} & \textbf{Role} \\
\midrule
SmallCNN & 398,662 & $128\times128$ & low-capacity control \\
ResNet18 & 11,196,006 & $224\times224$ & mid-capacity baseline \\
EfficientNet-B0 & 4,056,226 & $224\times224$ & efficient compound scaling \\
ConvNeXt-Tiny & 27,849,350 & $224\times224$ & modern conv block design \\
\bottomrule
\end{tabularx}
\end{table}

\section{Frozen protocol}

The protocol was fixed before the first full run and then held constant across
architectures and seeds, so that a difference between models describes the models.

\begin{table}[H]
\centering
\caption{Training protocol, identical across all reported runs}
\label{tab:protocol}
\tabfont
\begin{tabularx}{\textwidth}{@{}lY@{}}
\toprule
\textbf{Setting} & \textbf{Value} \\
\midrule
Optimiser & AdamW, $\beta = (0.9, 0.999)$ \\
Learning rate & $1\times10^{-4}$, cosine decay \\
Weight decay & $1\times10^{-4}$ \\
Batch size & 32 global, distributed over 2 GPUs \\
Epochs & 12 maximum, early stopping on validation macro-F1 \\
Precision & mixed (fp16) with dynamic loss scaling \\
Augmentation & random resized crop, horizontal flip, rotation $12^\circ$, colour jitter \\
Normalisation & ImageNet mean and standard deviation \\
Seed handling & one seed per run, seeds 42, 43, 44 \\
Checkpoint selection & best validation macro-F1, never test \\
\bottomrule
\end{tabularx}
\end{table}

An earlier protocol, referred to as protocol v1, could not reproduce the reference
result and is documented as a blocker in Chapter~\ref{ch:limitations}. All numbers in
this report come from the frozen protocol above.

\section{Evaluation design}

\subsection{Unit of independence}

The unit is the training seed, not the image. Images within a run are not
independent observations of model behaviour, because they share a fitted model. A
confidence interval computed over 236 images describes sampling variation in the
test set, not variation in training. Both are reported, and they are never
substituted for one another.

\subsection{Metrics}

Macro-F1 is the headline metric because the class distribution is uneven and
accuracy hides per-class collapse. Balanced accuracy is also reported, with the
explicit note that it equals macro recall by construction and therefore adds no
independent information. Per-class precision, recall, F1 and support appear in the
appendix, and the confusion matrix is reported in full.

\subsection{Cross-dataset evaluation}

Field evaluation uses the checkpoint selected on PlantVillage validation. No
PlantDoc image is used for training, model selection, threshold selection or
calibration. PlantDoc test is scored once per checkpoint.

\subsection{Intervention design}

To ask what the decision depends on, we swap content rather than inspect heat maps.
Each image receives a donor of a different class twice: once its leaf, once its
background. Adoption means the prediction moves to the donor's class. A same-class
donor provides the control: any movement there measures the disruption caused by the
edit itself rather than by the swapped content. Chance adoption is $1/38$, about
2.6\,\%.

\subsection{Statistical treatment}

Paired comparisons use a bootstrap over images where the comparison is within a
single checkpoint, and the paired $t$ interval over seeds where the comparison is
between training runs. With three seeds the latter has two degrees of freedom and is
correspondingly wide. It is reported as it is rather than as significance.

\section{Reproducibility}

All training ran on Kaggle with two Tesla T4 GPUs. The notebook is generated by a
builder script that embeds the repository sources byte for byte, so a hand-edited
cell cannot drift from the code under test. Random seeds, dependency versions and
commands are recorded per run. Local hardware is used only for data preparation,
evaluation of returned checkpoints, and document generation.
